\def\SigleNumero{STT-1900}
\def\NomCours{Méthodes statistiques pour l'ingénierie}
\def\DateEvaluation{Examen d'entraînement}
\def\HeureEvaluation{}
\def\Session{}
\def\NomEvaluation{Examen pratique --- Modules 1 à 7}
\def\Ponderation{47}
% Ne rien mettre entre les accolades si l'enseignant (ou l'enseignante) est aussi le coordonnateur (ou la coordonnatrice)
\def\Coordonnatrice{}
\def\Coordonnateur{}

% Ne rien mettre entre les accolades s'il y a plus d'une section
\def\Enseignant{}
\def\Enseignante{}

% Ne rien mettre entre les accolades s'il s'agit d'un cours à section unique
% Mettre le nom de chacun des enseignants et enseignantes entre les accolades
% Une case à cocher sera générée pour chacune des sections
\def\SectionA{}
\def\SectionB{}
\def\SectionC{}
\def\SectionD{}

\newcommand{\affichepoints}{}     % ← AVEC points
% \renewcommand{\affichepoints}{on} % ← SANS points (non vide)

\def\Directives{
\begin{enumerate}
\item Matériel autorisé:
\begin{itemize}
\item Calculatrice scientifique {\bf autorisée par la faculté des sciences et de génie};
\item Votre aide-mémoire d'une feuille recto-verso, écrit à la main.
\end{itemize}
\item Assurez-vous que cette évaluation comporte \numquestions ~exercices répartis sur \numpages{}~pages, pour un total de \numpoints{}~points.\
\item \textbf{Veuillez ne pas dégrafer votre examen.}
\item Vous devez\textbf{ justifier toutes vos réponses} par une démarche claire, sauf pour les questions à choix de réponse.\
\item Pour les questions à choix de réponse et les vrais ou faux, vous pouvez mettre un X, un crochet ($\checkmark$) ou remplir la case.\
\item Ce document est un examen d'entraînement (non officiel), construit à partir du style des examens intra du cours. Chaque partie d'une question est indépendante des autres.
\end{enumerate}
}

\providecommand{\Gradescope}{}


\documentclass{Evaluation}
\usepackage{multicol,graphicx}

% Laisser \noprintanswers en commentaire pour obtenir la version CORRIGÉE (avec solutions).
% Décommenter la ligne suivante pour obtenir la version ÉTUDIANTE (sans solutions).
% \noprintanswers

\begin{document}

\pageblanche
%%%%%%%%%%%%%%%% QUESTION %%%%%%%%%%%%%%%%
\begin{questions}

\question[4]\label{descriptive}
\noindent
\begin{minipage}[t]{0.5\linewidth}\vspace{0pt}
On mesure une variable continue chez 50 individus et on obtient la distribution de fréquences ci-contre.\\[0.5em]

\begin{tabular}{c|c}
Classe & Effectif \\ \hline
$[0,10[$ & 6 \\
$[10,20[$ & 14 \\
$[20,30[$ & 20 \\
$[30,40[$ & 8 \\
$[40,50[$ & 2
\end{tabular}
\end{minipage}\hfill
\begin{minipage}[t]{0.42\linewidth}\vspace{0pt}
\centering
\begin{tikzpicture}[x=0.2cm,y=0.35cm]
\draw[->, thick] (-1,0) -- (52,0);
\draw[->, thick] (0,0) -- (0,22);
\foreach \x in {0,10,20,30,40,50}
  \draw (\x,0) -- (\x,-0.5) node[below] {\x};
\foreach \y in {0,4,8,12,16,20}
  \draw (0,\y) -- (-1,\y) node[left] {\y};
\node at (25,-3) {Valeur};
\node[rotate=90] at (-6,11) {Effectif};
\draw (0,0) rectangle (10,6);
\draw (10,0) rectangle (20,14);
\draw (20,0) rectangle (30,20);
\draw (30,0) rectangle (40,8);
\draw (40,0) rectangle (50,2);
\end{tikzpicture}
\end{minipage}

\begin{parts}
\part[2] Dans quelle classe se situe le quantile d'ordre $0{,}60$? Justifiez à l'aide des effectifs cumulés.

\begin{solution}
Le rang visé est $0{,}60\times 50 = 30$. Les effectifs cumulés sont: $6$ (après $[0,10[$), $6+14=20$ (après $[10,20[$), $20+20=40$ (après $[20,30[$). Puisque $20 < 30 \le 40$, le quantile se situe dans la classe $[20,30[$.
\end{solution}

\part[2] Sachant que ce quantile se situe dans la classe $[20,30[$, calculez sa valeur exacte (par interpolation linéaire).

\begin{solution}
Avec $L=20$, effectif cumulé avant la classe $=20$, effectif de la classe $=20$, largeur $=10$:
$$q_{0{,}60} = 20 + \frac{30-20}{20}\times 10 = 20+5 = 25.$$
\end{solution}
\end{parts}

\vspace{0.3cm}

\question[4]\label{fiabilite} On considère deux montages électroniques indépendants l'un de l'autre.

\begin{parts}
\part[2] Un système est composé de deux composants $C_1$ et $C_2$ montés \textbf{en série} (le système fonctionne seulement si les deux composants fonctionnent). Les composants fonctionnent de façon indépendante, avec $P(C_1 \text{ fonctionne})=0,95$ et $P(C_2\text{ fonctionne})=0,85$. Quelle est la probabilité que le système ne fonctionne pas?

\begin{solution}
$P(\text{fonctionne}) = 0,95\times 0,85 = 0,8075$, donc $P(\text{ne fonctionne pas}) = 1-0,8075 = 0,1925$.
\end{solution}

\part[2] Un bloc est composé de trois composants $D_1,D_2,D_3$ montés \textbf{en parallèle} (le bloc fonctionne si au moins un composant fonctionne). Les composants fonctionnent indépendamment, chacun avec probabilité $0,6$. Quelle est la probabilité que le bloc fonctionne?

\begin{solution}
$P(\text{ne fonctionne pas}) = (1-0,6)^3 = 0,4^3 = 0,064$, donc $P(\text{fonctionne}) = 1-0,064 = 0,936$.
\end{solution}
\end{parts}

\pageblanche

\question[5]\label{densite} Soit $X$ une variable aléatoire continue de densité
$$f(x) = \frac{3}{8}x^2, \qquad 0 < x < 2, \qquad f(x)=0 \text{ ailleurs.}$$

\begin{parts}
\part[2] Vérifiez que $f$ est bien une fonction de densité valide.

\begin{solution}
$f(x)\ge 0$ sur $(0,2)$ et
$$\int_0^2 \frac{3}{8}x^2\,dx = \frac{3}{8}\left[\frac{x^3}{3}\right]_0^2 = \frac{3}{8}\times\frac{8}{3}=1,$$
donc $f$ est une densité valide.
\end{solution}

\part[2] Calculez $P(1<X<2)$.

\begin{solution}
$$P(1<X<2)=\int_1^2\frac{3}{8}x^2\,dx = \left[\frac{x^3}{8}\right]_1^2 = 1-\frac{1}{8}=\frac{7}{8}.$$
\end{solution}

\part[1] Calculez $E(X)$.

\begin{solution}
$$E(X)=\int_0^2 x\cdot\frac{3}{8}x^2\,dx = \frac{3}{8}\left[\frac{x^4}{4}\right]_0^2 = \frac{3}{8}\times 4 = 1,5.$$
\end{solution}
\end{parts}

\vspace{0.3cm}

\question[5]\label{normale} Soit $X\sim N(80,100)$.

\begin{parts}
\part[2] Calculez $P(X>95)$.

\begin{solution}
$Z=\dfrac{95-80}{10}=1,5$, donc $P(X>95)=P(Z>1,5)=1-0,9332=0,0668$.
\end{solution}

\part[1] Trouvez $x_0$ tel que $P(X<x_0)=0,10$.

\begin{solution}
On cherche $z$ tel que $P(Z<z)=0,10$, soit $z=-1,28$. Donc $x_0 = 80+(-1,28)\times 10 = 67,2$.
\end{solution}

\part[2] Soit $Y\sim N(50,64)$, indépendante de $X$. Calculez $P(X-Y>40)$.

\begin{solution}
$X-Y \sim N(80-50,\,100+64) = N(30,164)$. $Z=\dfrac{40-30}{\sqrt{164}}=\dfrac{10}{12,806}=0,78$, donc $P(X-Y>40)=P(Z>0,78)=1-0,7823=0,2177$.
\end{solution}
\end{parts}

\pageblanche

\question[6]\label{conjointeCombi}

\begin{parts}
\part[3] Soit $X$ et $Y$ deux variables aléatoires discrètes dont la loi conjointe est donnée par le tableau suivant:
$$
\begin{array}{c|cc}
 & X=0 & X=1 \\ \hline
Y=0 & 0,1 & 0,3 \\
Y=1 & 0,4 & 0,2
\end{array}
$$
Calculez $Cov(X,Y)$.

\begin{solution}
Lois marginales: $P(X=0)=0,1+0,4=0,5$, $P(X=1)=0,3+0,2=0,5$; $P(Y=0)=0,1+0,3=0,4$, $P(Y=1)=0,4+0,2=0,6$.

$$E(X)=0(0,5)+1(0,5)=0,5, \qquad E(Y)=0(0,4)+1(0,6)=0,6.$$

Seul le terme $(X,Y)=(1,1)$ contribue à $E(XY)$: $E(XY)=1\times1\times0,2=0,2$.

$$Cov(X,Y)=E(XY)-E(X)E(Y)=0,2-(0,5)(0,6)=0,2-0,3=-0,1.$$
\end{solution}
\vspace{4cm}

\part[3] Soit $U$, $V$ et $W$ trois variables aléatoires avec $E(U)=2$, $E(V)=-1$, $E(W)=3$; $V(U)=4$, $V(V)=9$, $V(W)=1$; $Cov(U,V)=-1$, $Cov(U,W)=0,5$, et $V$ indépendante de $W$. Calculez $E(2U-3V+W)$ et $V(2U-3V+W)$.

\begin{solution}
$$E(2U-3V+W)=2(2)-3(-1)+3 = 4+3+3=10.$$
$$V(2U-3V+W)=2^2V(U)+3^2V(V)+V(W)+2\big[2(-3)Cov(U,V)+2(1)Cov(U,W)+(-3)(1)Cov(V,W)\big]$$
$$= 4(4)+9(9)+1+2\big[(-6)(-1)+2(0,5)+0\big] = 16+81+1+2(6+1)=98+14=112.$$
\end{solution}
\end{parts}

\pageblanche

\question[6]\label{tclStudent}

\begin{parts}
\part[3] Soit $X_1,\ldots,X_{100}$ i.i.d. de loi quelconque avec $E(X_i)=20$ et $V(X_i)=64$. À l'aide du théorème central limite, calculez approximativement $P(\overline{X}>21,4)$.

\begin{solution}
Par le TCL, $\overline{X} \stackrel{\cdot}{\sim} N\!\left(20,\ \dfrac{64}{100}\right)=N(20,\,0,64)$, donc $\sigma_{\overline{X}}=0,8$.
$$Z=\frac{21,4-20}{0,8}=1,75 \quad\Rightarrow\quad P(\overline{X}>21,4)=P(Z>1,75)=1-0,9599=0,0401.$$
\end{solution}
\vspace{5cm}

\part[3] Soit $Y_1,\ldots,Y_{16}$ i.i.d. de loi $N(\mu,\sigma^2)$, et soit $S^2$ la variance échantillonnale. On donne $t_{0,025,15}=2,131$. Calculez
$$P\left(-2,131 < \frac{\overline{Y}-\mu}{S/\sqrt{16}} < 2,131\right).$$

\begin{solution}
La quantité $\dfrac{\overline{Y}-\mu}{S/\sqrt{16}}$ suit exactement une loi de Student à $16-1=15$ degrés de liberté. Donc
$$P(-2,131<t_{15}<2,131) = 1-2(0,025)=0,95.$$
\end{solution}
\end{parts}

\pageblanche

\question[7]\label{estimateurs} Soit $X_1,X_2,X_3,X_4$ un échantillon aléatoire simple provenant d'une variable d'espérance $\mu$ et de variance $\sigma^2$. On considère deux estimateurs de $\mu$:
$$\hat\theta_1=\frac{X_1+X_2+X_3+X_4}{4}, \qquad \hat\theta_2=\frac{X_1+2X_2+2X_3+X_4}{6}.$$

\begin{parts}
\part[4] Calculez $V(\hat\theta_1)$ et $V(\hat\theta_2)$ en fonction de $\sigma^2$, puis déterminez lequel des deux estimateurs a la plus petite variance.

\begin{solution}
$$V(\hat\theta_1)=\frac{4\sigma^2}{16}=\frac{\sigma^2}{4}=0,25\,\sigma^2.$$
$$V(\hat\theta_2)=\frac{1^2+2^2+2^2+1^2}{6^2}\sigma^2=\frac{10}{36}\sigma^2=\frac{5}{18}\sigma^2\approx 0,2778\,\sigma^2.$$
Puisque $0,25\sigma^2 < 0,2778\sigma^2$, $\hat\theta_1$ a la plus petite variance.
\end{solution}
\vspace{6cm}

\part[3] On sait que $biais(\hat\theta_2,\mu)=0,1$ (une constante) et que $V(\hat\theta_2)=\dfrac{5\sigma^2}{18}$. En supposant $\sigma^2=9$, calculez l'erreur quadratique moyenne (EQM/MSE) de $\hat\theta_2$.

\begin{solution}
$$MSE(\hat\theta_2,\mu) = biais(\hat\theta_2,\mu)^2 + V(\hat\theta_2) = 0,1^2 + \frac{5}{18}\times 9 = 0,01+2,5=2,51.$$
\end{solution}
\end{parts}

\pageblanche

\question[6]\label{intervalles}

\begin{parts}
\part[3] Soit $X_1,\ldots,X_{49}$ i.i.d. $N(\mu,\sigma^2=196)$, avec $\overline{x}=85$. Construisez un intervalle de confiance à 95\% pour $\mu$.

\begin{solution}
$\sigma=14$, donc $\sigma/\sqrt{49}=14/7=2$. Avec $z_{0,025}=1,96$:
$$IC_{95\%} = \left[85-1,96(2),\ 85+1,96(2)\right] = [81,08,\ 88,92].$$
\end{solution}
\vspace{5cm}

\part[3] Soit $Y_1,\ldots,Y_{10}$ i.i.d. $N(\mu,\sigma^2)$ (variance inconnue), avec $\overline{y}=42$ et $s=6$. On donne $t_{0,05,9}=1,833$. Construisez un intervalle de confiance à 90\% pour $\mu$.

\begin{solution}
$s/\sqrt{10}=6/3,1623=1,897$. Avec $t_{0,05,9}=1,833$:
$$IC_{90\%}=\left[42-1,833(1,897),\ 42+1,833(1,897)\right]=[42-3,48,\ 42+3,48]=[38,52,\ 45,48].$$
\end{solution}
\end{parts}

\pageblanche

\question\label{echantillonnage} Soit $K\sim\chi^2_6$, $T\sim t_{10}$ et $F\sim F_{4,8}$, trois variables aléatoires indépendantes. Dites si chaque énoncé est vrai ou faux.

\begin{parts}
\part[1] $E(K)=6$.
\begin{oneparcheckboxes}
\correctchoice Vrai
\choice Faux
\end{oneparcheckboxes}

\begin{solution}
Vrai. L'espérance d'une $\chi^2_n$ est $n$; ici $n=6$.
\end{solution}

\part[1] $V(T)=1,25$.
\begin{oneparcheckboxes}
\correctchoice Vrai
\choice Faux
\end{oneparcheckboxes}

\begin{solution}
Vrai. $V(t_n)=\dfrac{n}{n-2}$ pour $n>2$; ici $\dfrac{10}{8}=1,25$.
\end{solution}

\part[1] $1/F$ suit une loi de Fisher à 8 et 4 degrés de liberté.
\begin{oneparcheckboxes}
\correctchoice Vrai
\choice Faux
\end{oneparcheckboxes}

\begin{solution}
Vrai. Propriété de la loi de Fisher: si $F\sim F_{n,m}$, alors $1/F\sim F_{m,n}$; ici $1/F\sim F_{8,4}$.
\end{solution}

\part[1] La loi de Student est asymétrique.
\begin{oneparcheckboxes}
\choice Vrai
\correctchoice Faux
\end{oneparcheckboxes}

\begin{solution}
Faux. La loi de Student est symétrique autour de 0 (comme la loi normale).
\end{solution}
\end{parts}

\end{questions}
\pageblanche

\end{document}